\documentclass{article}

\usepackage{corl_2026} % Use this for the initial submission.
% \usepackage[final]{corl_2026} % Uncomment for the camera-ready ``final'' version.
% \usepackage[preprint]{corl_2026} % Uncomment for pre-prints (e.g., arxiv); This is like ``final'', but will remove the CORL footnote.

\usepackage{url}
% \usepackage{refcheck}
% \usepackage[backref]{hyperref} %
\usepackage{cite}
\usepackage{amsmath,amssymb,amsfonts}
\usepackage{graphicx}
\usepackage{mathtools}
% \usepackage[ruled]{algorithm2e}
\usepackage[T1]{fontenc}
\usepackage[dvipsnames]{xcolor}
\usepackage{subfig}
\usepackage{multirow}
\usepackage{makecell} % for \Xhline
\usepackage{optidef} % For typesetting optimization problems
\usepackage{siunitx}
\sisetup{
    round-mode          = places,
    round-precision     = 4,
    table-number-alignment = center,
    detect-weight       = true,
    detect-inline-weight = math
}
\allowdisplaybreaks

% Formatting and General
\newcommand{\etal}{\textit{et al. }}
\newcommand{\solution}{\textbf{Solution:}}
\newcommand{\qed}{\hfill \ensuremath{\Box}}
\definecolor{darkgreen}{rgb}{0.0, 0.5, 0.0} 

% Math Sets and Symbols
\newcommand{\reals}{\ensuremath{\mathbb{R}}}
\newcommand{\naturals}{\ensuremath{\mathbb{N}}}
\newcommand{\integers}{\ensuremath{\mathbb{Z}}}
\newcommand{\expectation}{\ensuremath{\mathbb{E}}}
\newcommand{\D}{\mathrm{d}}
\newcommand{\reg}{\ensuremath{\emph{R}}}

% Spaces and Sets
\newcommand{\config}{\mathcal{Q}}
\newcommand{\wspace}{\mathcal{W}}
\newcommand{\obst}{\mathcal{O}}
\newcommand{\goals}{\mathcal{G}}
\newcommand{\states}{\mathcal{S}}
\newcommand{\trans}{\ensuremath{\mathcal{T}}}

% Rewards
\newcommand{\rterm}{R_{\mathrm{term}}}
\newcommand{\rspray}{R_{\mathrm{spr}}}
\newcommand{\rpath}{R_{\mathrm{path}}}
\newcommand{\rvisit}{R_{\mathrm{visit}}}

% Theorem Environments
\newtheorem{problem}{Problem}
\newtheorem{remark}{Remark}
\newtheorem{definition}{Definition}
\newtheorem{proposition}{Proposition}
\newtheorem{lemma}{Lemma}
\newtheorem{theorem}{Theorem}

% Paired Delimiters
\DeclarePairedDelimiter\ceil{\lceil}{\rceil}
\DeclarePairedDelimiter\floor{\lfloor}{\rfloor}

% Operators
\DeclareMathOperator*{\argmax}{arg\,max}
\DeclareMathOperator{\simulate}{SIM}
\DeclareMathOperator{\init}{init}
\DeclareMathOperator{\free}{free}
\DeclareMathOperator{\new}{new}
\DeclareMathOperator{\pos}{POS}
\DeclareMathOperator{\goalf}{GS}
\DeclareMathOperator{\visited}{\mathbb{V}}

% Fallback definitions for general-formulation macros (no-op if already defined in macros.tex)
\providecommand{\obst}{\mathcal{O}}
\providecommand{\free}{\mathrm{free}}
\providecommand{\obs}{\mathrm{obs}}
\providecommand{\col}{\mathrm{col}}

\title{Learning to Navigate in a Stochastic and Uncertain Environment with Different Robot Dynamics}

% The \author macro works with any number of authors. There are two
% commands used to separate the names and addresses of multiple
% authors: \And and \AND.
%
% Using \And between authors leaves it to LaTeX to determine where to
% break the lines. Using \AND forces a line break at that point. So,
% if LaTeX puts 3 of 4 authors names on the first line, and the last
% on the second line, try using \AND instead of \And before the third
% author name.

% NOTE: authors will be visible only in the camera-ready and preprint versions (i.e., when using the option 'final' or 'preprint'). 
% 	For the initial submission the authors will be anonymized.

\author{
  Jane E.~Doe\\
  Department of Electrical Engineering and Computer Sciences\\
  University of California Berkeley 
  United States\\
  \texttt{janedoe@berkeley.edu} \\
  %% examples of more authors
  %% \And
  %% Coauthor \\
  %% Affiliation \\
  %% Address \\
  %% \texttt{email} \\
  %% \AND
  %% Coauthor \\
  %% Affiliation \\
  %% Address \\
  %% \texttt{email} \\
  %% \And
  %% Coauthor \\
  %% Affiliation \\
  %% Address \\
  %% \texttt{email} \\
  %% \And
  %% Coauthor \\
  %% Affiliation \\
  %% Address \\
  %% \texttt{email} \\
}


\begin{document}
\maketitle

%===============================================================================

\begin{abstract}
    Multi-robot navigation in real-world settings is inherently \emph{stochastic} and \emph{partially predictable}: wind perturbs motion, actuators do not match nominal commands, and onboard sensing is noisy. Yet many reinforcement learning (RL) navigation benchmarks either omit these effects or treat them informally, making it hard to compare algorithms under realistic uncertainty and across heterogeneous platforms. This paper studies learning-based multi-robot navigation under such disturbances while explicitly contrasting \emph{two} continuous-control dynamics families: planar omnidirectional aerial vehicles modeled as force-driven point masses, and nonholonomic wheeled ground robots modeled with bicycle kinematics in cluttered workspaces. We retain a dynamics-agnostic constrained-optimization view of ROI coverage with collision avoidance, derive a corresponding finite-horizon Markov decision process (MDP) reward decomposition, and summarize an equivalence-style argument linking the navigation objective to optimal MDP behavior under idealized full-state observability. We then describe an open, Gymnasium-compatible implementation that augments this idealization with stochastic wind, action noise, observation noise, optional domain randomization, and ablation switches for rewards, observations, and uncertainty modes. Using Stable-Baselines3 and SB3-Contrib, we benchmark six standard deep RL algorithms with default and Optuna-tuned hyperparameters across \emph{ten} randomized aerial field layouts and \emph{eight} wheeled layouts. Finally, we visualize aerial policies in a CoppeliaSim agricultural scene and analyze wind-strength sensitivity; a companion line of work studies selective \emph{spraying} under richer deposition and belief-state structure.
\end{abstract}

% Two or three meaningful keywords should be added here
\keywords{CoRL, Robots, Learning} 

%===============================================================================

\section{Introduction}
\label{sec:intro}
Developing robots that can plan and navigate in complex, uncertain environments is a fundamental problem in artificial intelligence. Reinforcement learning (RL) is a widely used paradigm for such tasks \citep{chen2017icra, everett2018iros, long2017iral, fan2020ijrr, han2022iral}. Progress in multi-robot continuous control, however, depends not only on algorithms but also on \emph{environments} that expose stochasticity, sensor limitations, and heterogeneous dynamics in reproducible ways \citep{brockman2016openai, islam2017reproducibility, mania2018ars, dulacArnold2021, sutton2018reinforcement}. Once formulated as Markov (or partially observable Markov) decision processes, such environments become shared testbeds for gradient-based RL and for derivative-free alternatives such as Augmented Random Search (ARS) \citep{mania2018ars}, evolution strategies \citep{salimans2017}, and other black-box optimization techniques \citep{risi2015neuroevolution}. Unfortunately, many RL benchmark papers omit a complete MDP description \citep{kurach2020google, zhang2019cityflow} or provide only informal reward engineering for navigation \citep{chen2017icra, everett2018iros, long2017iral, fan2020ijrr, han2022iral}, and they seldom evaluate the \emph{same} training stack under controlled uncertainty factors across \emph{different} robot models.

In this paper, we center the question: \emph{how do standard deep RL methods behave when multi-robot navigation is deliberately made stochastic and uncertain, and when the robot dynamics differ substantially between platforms?} Concretely, we pair (i) a principled optimization-to-MDP construction for ROI visitation with collision avoidance, with (ii) a public codebase that instantiates \emph{two} dynamics classes---aerial point-mass UAVs and wheeled bicycle robots---under configurable wind volatility, actuator noise, observation noise, and domain randomization, all behind a unified training interface. We benchmark common gradient-based algorithms (A2C \citep{mnih2016a2c}, TRPO \citep{schulman2015trust}, PPO \citep{schulman2017ppo}, TQC \citep{kuznetsov2020tqc}, CrossQ \citep{bhatt2024crossq}) together with ARS \citep{mania2018ars}, comparing Stable-Baselines3 defaults against Optuna-selected hyperparameters \citep{optuna_2019, henderson2018deep}. Training and evaluation use Gymnasium \citep{towers2024gymnasium} with two-dimensional Pygame visualization for \emph{both} platforms; for UAVs we additionally export policies to a CoppeliaSim scene for qualitative inspection in an agricultural-style setting. Importantly, the Gymnasium task studied here is \emph{target visitation / coverage} (binary ROI completion), whereas our parallel manuscript treats cooperative \emph{precision spraying} with deposition uncertainty and a partially observable formalization.

While multi-robot navigation has previously been formulated within RL, few works simultaneously (a) articulate an optimization-aligned MDP specification, (b) stress-test identical algorithmic baselines under multi-source uncertainty, and (c) keep the same experimental switches across qualitatively different dynamics. Relative to our earlier emphasis on ``environment design'' as the headline, the present title highlights \emph{learning under stochasticity} and \emph{cross-dynamics} evaluation. To summarize, the contributions are:
\begin{itemize}
    \item \textbf{Formal design template:} We present a finite-horizon MDP reward decomposition for multi-robot ROI coverage with collision penalties and redundant-motion costs, together with an analytical linkage to a constrained navigation optimization problem under idealized assumptions (Section~\ref{sec:problem}--\ref{sec:mdp_theory}).
    \item \textbf{Stochastic, uncertainty-aware benchmark:} We release Gymnasium environments with stochastic wind, action noise, observation noise, optional domain randomization, and ablations over reward terms, observation channels, and uncertainty modes, enabling controlled studies of robustness (Section~\ref{sec:experiments}).
    \item \textbf{Cross-dynamics evaluation:} Using one training codebase, we evaluate six deep RL algorithms on \emph{ten} randomized aerial fields and \emph{eight} wheeled obstacle layouts, reporting episodic return statistics under default and tuned hyperparameters.
    \item \textbf{Simulation and transfer-oriented analysis:} We visualize aerial behaviors in CoppeliaSim, perform wind-strength sensitivity at inference time, and discuss observation-noise stress tests; a separate companion paper focuses on selective spraying with richer physical models.
\end{itemize}

\section{Motivation and Related Work}
\label{sec:related}
Over the past two decades, numerous methods have been developed for multi-robot navigation and control \citep{nafis2021, sam2005, mickey2021}. Classical planning approaches such as A* \citep{goto2003}, RRT \citep{rodriguez2006rrt}, RRT* \citep{karaman2011rrts}, PSO \citep{nafis2021}, and RA-MCPP \citep{mickey2021} rely on explicit modeling of path-planning heuristics or global optimization strategies. While effective in structured settings, these methods often become computationally expensive in high-dimensional continuous state and action spaces, particularly when multiple robots and dynamic environments are involved. As a result, scalable multi-robot navigation in complex environments remains a fundamental challenge in robotics \citep{Sunderhauf2018ijrr, ibarz2021ijrr, zhang2023multi, cai2022multi, osooli2024multi}.

Reinforcement Learning (RL) has emerged as a promising alternative, especially for problems requiring real-time decision-making and adaptation in high-dimensional spaces \citep{kober2014, chen2017icra, everett2018iros, bae2019as, tang2025, dulacArnold2021, ibarz2021ijrr}. Early RL-based navigation works \citep{kober2014, bae2019as} often considered discrete or simplified state-action spaces. More recent approaches \citep{chen2017icra, everett2018iros, long2017iral, fan2020ijrr, han2022iral} formulate navigation objectives as optimization problems and derive corresponding Markov Decision Processes (MDPs). However, these formulations typically lack formal guarantees establishing equivalence between the navigation objective and the optimal policy. In addition, many works assume fixed robot dynamics, rely heavily on reward heuristics, or introduce customized RL architectures tailored to specific environments. Beyond navigation, several benchmark environments have been proposed to bridge RL and real-world challenges. Dulac-Arnold \textit{et al.}~\citep{dulacArnold2021} introduced a suite of practical RL challenges and corresponding MDP formulations, but none address multi-robot navigation, and formal correctness analyses remain limited. Large-scale environments such as Google Research Football \citep{kurach2020google} and CityFlow \citep{zhang2019cityflow} provide realistic testbeds but do not offer formal guarantees relating task objectives to optimal policies. Furthermore, the need for reproducible and extensible baselines in continuous control has been emphasized by Mania \textit{et al.}~\citep{mania2018ars}, who showed that algorithmic performance can vary significantly in high-dimensional settings.

In contrast to prior works, we focus on the principled construction of the RL environment itself \emph{and} on exposing stochasticity in a tunable, reproducible way. We provide a complete MDP specification for multi-robot navigation and formally prove that, under a finite-horizon setting and appropriate reward design, the optimal MDP policy is equivalent to the solution of the underlying constrained optimization problem in the idealized setting of Section~\ref{sec:mdp_theory}. The accompanying software extends this idealization with stochastic disturbances and noisy observations for empirical RL studies. Moreover, our framework accommodates arbitrary robot dynamics through the generic simulator map $\simulate(\cdot)$ and can be solved using standard state-of-the-art RL algorithms without task-specific architectural modifications. By isolating the contribution to environment design and validating it across diverse dynamics and algorithms, we position our framework as both theoretically grounded and practically extensible for high-dimensional multi-robot systems.

\section{Preliminaries}
\label{sec:background}
Let $\reals, \reals^+, \naturals, \integers$ be the set of reals, positive reals, naturals, and integers, respectively. For a set $S$, let $|S|$ and $\mathbb{P}(S)$ denote its length and power set, respectively.  Given two real numbers $x,y \in \reals$, let $[x,y] = \{ i \in \reals \mid x \leq i \leq y \}$ denote the set of all real numbers from x to y. Let $Y$ be a continuous random variable with probability density function $f$. The expected value (or expectation) is defined as, $\expectation[Y] = \int_{-\infty}^{\infty} y f(y)dy$. For a set of samples $S$, let $\Delta(S)$ denote the set of probability distributions over $S$. For a nonempty set $X$, a collection $\mathcal{A} \subseteq \mathcal{P}(X)$ is called $\sigma$-algebra \citep{tao2011} on $X$ if it satisfies $X \in \mathcal{A}$, $\forall A_i \in \mathcal{A}: X \setminus A_i \in \mathcal{A}$ and $ \cup_i^{\infty}A_i \in \mathcal{A}$. Let $(X, \mathcal{A})$ and $(Y, \mathcal{B})$ be measurable spaces where $X$ and $Y$ are sets and $\mathcal{A}$ and $\mathcal{B}$ are $\sigma$ algebras of subsets of $X$ and $Y$. A function $f: X \to Y$ is called measurable if $\forall B \in \mathcal{B}: f^{-1}(B)\in \mathcal{A}$. In RL, a policy $\pi: S \to \Delta (A)$ is a function that maps a state $s \in S$ to a probability distribution over the set of actions $A$. The action-value or Q-value of a policy $\pi$ at state $s$ and action $a$ is defined as the expected cumulative reward from following the policy after taking the action $a$ starting from $s$ and written as, $Q^{\pi}(s,a) = \expectation_{\pi}[\sum_{t\geq0}\gamma^t r_t | s_0=s, a_0=a]$ where $r_t \in \reals^+$ is the reward at time $t$ and $\gamma \in [0,1]$ is the discount factor. Following the notations and definitions from Choset \textit{et al.}\citep{choset2005}, we define the multi-robot navigation problem in the next section.

\section{Problem Formulation}
\label{sec:problem}
Consider a group of $n$ mobile robots $\mathcal{M} = \{ \mathcal{M}_1, \ldots, \mathcal{M}_n\}$ operating in a $d$-dimensional workspace $\wspace \subseteq \reals^d$ containing static obstacles $\obst$ and $m$ regions of interests (ROIs) $\goals = \{ \goals_1, \ldots, \goals_m\}$. Let $\wspace_{\obst_j}$ be the region occupied by obstacle $j$, and $\cup_j \wspace_{\obst_j}$ is the total region covered by all obstacles in the workspace. Therefore, the free workspace is the region $\wspace^{\free} = \wspace \setminus \cup_j \wspace_{\obst_j}$. From Plaku \textit{et al.} \citep{plaku2021}, each robot model $\mathcal{M}_i = \{ \mathcal{P}_i, \config_i, \states_i, s_i^{\init}, \mathcal{U}_i, f_i\}$ is defined in terms of its shape $\mathcal{P}_i$, configuration space $\config_i$, state space $\states_i$, initial state $s_i^{\init} \in \states_i$, control space $\mathcal{U}_i$ and dynamics $f_i$. Given the shape $\mathcal{P}_i$ for robot $i$, let $\mathcal{P}_i(q_i)\subset \wspace$ and $\mathcal{P}_i(s_i) \subset \wspace$ denote the regions occupied by the robot in the workspace $\wspace$ at configuration $q_i \in \config_i$ and state $s_i \in \states_i$ respectively. The configuration space occupied by all obstacles for the robot $i$ is, $\config_i^{obs} = \{ q_i \in \config_i \mid \exists j: \mathcal{P}_i(q_i) \cap \wspace_{\obst_j} \neq \emptyset \}$. The collision configuration space for robot $i$ is the set of configurations that collide with other robots in the workspace:
\begin{flalign} 
\begin{split}
\config_i^{col} = \{q_i \in \config_i \mid \exists j \neq i: (\mathcal{P}_i(q_i) \cap \mathcal{P}_j(q_j) \neq \emptyset) \} \cup \config_i^{obs}
\end{split}
\end{flalign}
Since the robots are moving, let $\config_i^{col}(t)$ denote the collision configuration for robot $i$ at time $t$. Therefore, the free configuration space for the robot $i$ at time $t$ is, $\config_i^{\free}(t) = \config_i \setminus \config_i^{col}(t)$. The state space $\states_i = \config_i \times \mathcal{V}_i$ of robot $i$ augments the configuration space $Q_i$ and motion space $\mathcal{V}_i$, where $\mathcal{V}_i$ contains the motion parameters such as steering angle, velocity, heading direction, etc. The kinematics and dynamics of each robot $i$ is described as follows:
\begin{align}
    \dot{q}_i &= f_i(q_i,u_i), q_i \in \config_i, u_i \in \mathcal{U}_i &\text{(kinematics)} &\label{eq:kinematics}\\
    \dot{s}_i &= f_i(s_i,u_i), s_i \ \in \states_i, u_i \in \mathcal{U}_i &\text{(dynamics)} &\label{eq:dynamics}
\end{align}
A new state $s_i^{\new} \in \states_i$ can be reached after a small time step $dt$ by numerically integrating $f_i$ by a function:
\begin{equation}
    s_i^{\new} \gets \simulate(s_i, u_i, f_i, dt)
\end{equation}
Let $\zeta_i: [0,T] \to \states_i$ denote a trajectory function for robot $i$ for time period $[0,T]$. Let $\goalf(\zeta_i(\tau))$ denote the ROIs visited by trajectory $\zeta_i$ from time $0$ to $\tau$, i.e.,
\begin{equation}
\begin{split}
\goalf(\zeta_i(\tau)) = \{ g \mid g \in \goals \land \exists t \in [0,\tau]: \mathcal{P}_i(\zeta_i(t)) \cap g \neq \emptyset \}
\end{split}
\end{equation}
We define the multi-robot navigation problem as a constrained optimization problem. Let us define a simple cost function $\mathcal{C}_i: \states_i \to \{ c_0, c_1 \}$ where $\mathcal{C}_i(s_i) = c_0$ if state $s_i$ is not visited before and $\mathcal{C}(s) = c_1$ if $s$ is visited before by robot $i$, and $c_0 < c_1$ (visited locations has higher costs). The objective is to generate $n$ trajectories $\zeta_1,\ldots.,\zeta_n$ for the robots by applying $n$ control signals $u_1,\ldots,u_n$ that solves the optimization problem:
\begin{align}
& \qquad \qquad \min_{\zeta_1, \ldots, \zeta_n} \bigg( \sum_{i=1}^{n} \int_{t=0}^{T} \mathcal{C}_i (\zeta_i(t))dt \bigg)   \label{eq:opt} \\
& \textrm{subject to, }    \nonumber \\
& (i) \quad \bigcup_{i=1}^n \goalf(\zeta_i(T)) = \goals  \label{eq:const1}\\
& (ii) \quad \forall i \in \{1,\ldots,n\}, \forall t \in [0,T]: \nonumber \\
& \qquad \zeta_i(t) = (q_i, v_i), q_i \in \config_i^{\free}(t), v_i \in \mathcal{V}_i \label{eq:const2}\\
& (iii) \quad \forall i \in \{1,\ldots,n\}, \forall t \in [0,T]: \nonumber \\
& \qquad \zeta_i(t+dt) \gets \simulate( \zeta_i(t), u_i, f_i, dt) \label{eq:const3}
\end{align}
Here, the first constraint ensures that all regions of interest or goals are visited. The second constraint ensures that there are no collisions. The third constraint ensures that each trajectory is dynamically feasible.

\section{Environment Design}
\label{sec:mdp_theory}
The above optimization problem can be represented as a reinforcement learning (RL) problem. We define a Markov Decision Process (MDP) $M = (S, A, \trans, R, \gamma)$ for the multi-robot navigation problem where:
• $S = \states_1 \times \ldots \times \states_n \times \{0, 1 \}^{|\goals|}$ is the set of states defined as a joint state space of $n$ robots, and a bitmask $b_t: \goals \to \{0, 1\}$ that encodes how many regions of interest are visited at time $t$. Therefore: 
\begin{equation}
\forall g \in \goals: b_t(g) =
\begin{cases}
    1, &\text{if $g$ is visited before $t$} \\
    0, &\text{otherwise}
\end{cases}
\end{equation}
• $A = \mathcal{U}_1 \times \ldots \times \mathcal{U}_n$ is the set of actions defined as the joint control inputs of $n$ robots
• $\trans: S \times A \to S$ is the transition function where at time $t$ for state $s_t$ and action $a_t$, it is defined as $\trans(s_t,a_t) = \Big( \simulate(s_1,u_1,f_1,dt), \ldots , \simulate(s_n,u_n,f_n,dt), b_{t+dt} \Big)$ where $s_t = (s_1,\ldots,s_n, b_t)$ , $a_t = (u_1,\ldots,u_n)$, and $f_1,\ldots, f_n$ are the dynamics of the $n$ robots. The $b_{t+dt}$ encodes all the goal positions reached by the robots at time $t+dt$.
• $R: S \times A \to \reals$ is the reward function, which is the sum of three different functions $R_C, R_G, R_V: S \times A \to \reals$ (designed for the three constraints of equation (\ref{eq:opt})) that at time $t$ satisfy:
\begin{align}
\label{eq:reward}
    R(s_t,a_t) = R_C(s_t,a_t) + R_G(s_t,a_t) + R_V (s_t,a_t)
\end{align}
Let $r_s, r_l \in \reals^+$ be some arbitrarily fixed small and large rewards where $r_s < r_l$. Let $r_M$ be a very large reward such that $r_M \gg nr_lT_{\max} $ where $T_{\max}$ is the maximum episode length. The first reward function $R_C$ determines the reward for avoiding collisions:
\begin{align}
\label{eq:reward_c}
& R_C(s_t, a_t) = \hfill \nonumber \\
& \begin{cases}
    -r_M, &\text{if } \exists i \in \{1,\ldots,n\}: s_i = (q_i, v_i) \land \\
          & q_i \notin \config_i^{\free}(t) \land v_i \in \mathcal{V}_i \\
    0,    &\text{otherwise}
\end{cases}
\end{align}
The second reward function $R_G$ specifies the reward the robots receive for visiting any region of interest (ROI). Let $\goals^{new}_{t+dt}$ be the newly visited ROIs by taking an action $a_t$ from state $s_t$:
\begin{equation}
    \goals^{new}_{t+dt} = \{ g \in \goals \mid b_t(g) = 0 \land b_{t+dt}(g) = 1 \}
\end{equation}
Therefore:
\begin{align}
R_G(s_t, a_t) = 
\begin{cases}
    +r_M, &\text{if } \forall g \in \goals: b_{t+dt} (g) = 1 \\
    r_l \cdot |\goals^{new}_{t+dt}|, & \text{otherwise} 
\end{cases}
\end{align}
The third reward function $R_V$ determines the reward that the robots receive by checking whether the locations in the state $s_t$ were visited previously in the episode. Let $V: {\wspace} \to \{0,1\}$ be a shared occupancy map where $V(x) = 1$ if the location $x \in \wspace$ is visited before $t$ in the current episode and 0 otherwise. We define the number of locations in the workspace that were visited by the $n$ robots before $t$ in the current episode as:
\begin{align}
n_{v_t} = \sum_{i=1}^n \{ i \mid \exists x \in \mathcal{P}_i(s_i): V(x) = 1 \}
\end{align}
Therefore:
\begin{align}
R_V(s_t, a_t) = - r_l n_{v_t} - r_s (n - n_{v_t}) \label{eq:reward_V}
\end{align}
• $\gamma \in [0,1]$ is the discount factor that is applied to future rewards.
For an infinite time horizon ($T=\infty$), the discount factor $\gamma$ must be less than 1 to converge $\sum_{t\geq0}\gamma^t r_t$ to a finite value when measuring the Q-value of a state and action. But for a finite time horizon ($T < \infty$), $\gamma = 1$ means that the reward for each timestep doesn't get discounted, and the cumulative reward is still finite and valid. This is a valid assumption since robots have power constraints and can only run for a limited time. Let $T_1, \ldots, T_n$ be the maximum runtimes of the $n$ robots. Therefore, we take the maximum length of each finite episode to be $T_{\max} = \min(T_1,\ldots,T_n)$. As discussed by Sutton and Barto \citep{sutton2018reinforcement}, the return objective takes future rewards into account more strongly as $\gamma$ approaches 1. Based on this discussion, we define the following theorem:
\begin{theorem}
    For a finite time horizon $T$, the optimal policy $\pi^*$ for the MDP $M$ will generate the paths that solve the optimization problem described in equation (\ref{eq:opt}) if $\gamma = 1$.
\end{theorem}
\textbf{Proof:} The definition of optimal policy $\pi^*$ and optimal action-value function $Q^{\pi^*}$ \citep{sutton2018reinforcement} are as follows:
\begin{flalign}
& \phantom{=\ } \pi^*(s_t) = \argmax_{a}Q^{\pi^*}(s_t,a_t) \\
& \phantom{=\ } Q^{\pi^*}(s_t,a_t) = \max_{\pi}Q^{\pi}(s_t,a_t) \\
& = \max_{\pi} \mathbb{E} [\sum_{t \geq 0} \gamma^t r_t \mid s_0=s_t, a_0=a_t, \pi] \nonumber \\
& = \max_{\pi} \mathbb{E} [\sum_{t \geq 0} \gamma^t R(s_t,a_t)  \mid s_0=s_t, a_0=a_t, \pi] \nonumber \\
& \rightline{[$r_t = R(s_t,a_t)$]} \nonumber \\ 
& = \max_{\pi} \mathbb{E} [\sum_{t \geq 0} \gamma^t (R_C(s_t,a_t) + R_G(s_t,a_t) +  R_V (s_t,a_t)) \nonumber \\ 
& \rightline{$ \mid s_0=s_t, a_0=a_t, \pi]$ [From equation (\ref{eq:reward})]} \nonumber 
\end{flalign}
From equation (\ref{eq:reward_c}), the maximum value of $R_{C}(s_t,a_t)$ is zero, which by definition means there are no collisions between any two robots or a robot and an obstacle, satisfying the second constraint of equation (\ref{eq:opt}). The maximum value of $R_G(s_t,a_t)$ is $r_M$ (since $r_M \gg nr_lT_{max}$), which means all the regions of interest are visited, satisfying the first constraint of equation (\ref{eq:opt}). From the definition of the transition function:
\begin{flalign}
&\trans ( s_t, a_t)  \nonumber \\
&= (\simulate(s_1,u_1,f_1,dt), \ldots , \simulate(s_n,u_n,f_n,dt), b_{t+dt}) \nonumber \\
&= (\simulate(\zeta_1(t),u_1,f_1,dt), \ldots , \simulate(\zeta_n(t), u_n,f_n, dt), b_{t+dt}) \nonumber \\
&= (\zeta_1(t+dt), \ldots , \zeta_n(t+dt), b_{t+dt})
\end{flalign}
This satisfies the third constraint of equation (\ref{eq:opt}). For a finite time horizon $T$ (each episode), we have the optimal action-value function:
\begin{flalign}
    & \phantom{=\ } Q^{\pi^*}(s_t,a_t) = \max_{\pi} \mathbb{E} [\sum_{t=0}^T \gamma^t (R_C(s_t,a_t) +  R_G(s_t,a_t) + \nonumber \\
    & \rightline{$ R_V(s_t,a_t)) \mid . ]$} \nonumber \\
    &= 0 + r_M + \max_{\pi} \mathbb{E} [\sum_{t=0}^T \gamma^t ( R_V(s_t,a_t)) \mid . ] \nonumber \\
    & \rightline{[From the above discussion]} \nonumber \\
    &= r_M + \max_{\pi} \mathbb{E} [\sum_{t=0}^T ( - r_l n_{v_t} - r_s (n - n_{v_t}) \mid . ] \nonumber \\
    & \rightline{[From equation (\ref{eq:reward_V}) and $\gamma^t=1$]} \nonumber \\
    &= r_M + \min_{\pi} \mathbb{E} [  \sum_{t=0}^T (r_l n_{v_t} +  r_s (n - n_{v_t}) \mid . ] \\
    & \rightline{[$\max\limits_{a\in A}(-f(a)) = \min\limits_{a\in A}(f(a))]$} \nonumber
\end{flalign}
From equation (\ref{eq:opt}), let $\mathcal{C}_i(\zeta_i(t)) = \mathcal{J}_i(t)$ be the cost for robot $i$ at time $t$. Since the robots can only be in one state $s \in \states_i$ at a time $t$, $\zeta^{-1}(s)$ will return the time that the robot was at that state if it was visited, and false otherwise. $\mathcal{C}^{-1}(c_0)$ and $\mathcal{C}^{-1}(c_1)$ will return the set of states that has the costs $c_0$ and $c_1$ respectively, and each set is a subset of $\states_i$. Therefore, $\mathcal{J}_i(t)$ or $\mathcal{C}_i(\zeta_i(t))$ is a measurable function. From the theorems of Fubini and Tonelli \citep{Royden2010RealAnalysis}, we have $\sum_a \int_b f(a,b)db = \int_b \sum_a f(a,b)db $ if the total sum is finite and $f(a,b)$ is measurable. Therefore from equation (\ref{eq:opt}):
\begin{align}
& \min \bigg( \sum_{i=1}^{n} \int_{t=0}^{T} \mathcal{C}_i (\zeta_i(t))dt \bigg) \nonumber = \min \bigg( \int_{t=0}^{T}  \sum_{i=1}^{n} \mathcal{C}_i (\zeta_i(t))dt \bigg) \nonumber \\
&=  \min \bigg( \int_{t=0}^{T} (\sum_{i=1}^{n_{v_t}} \mathcal{C}_i (\zeta_i(t)) + \sum_{i=n_{v_t}+1}^{n} \mathcal{C}_i (\zeta_i(t))) dt \bigg) \nonumber \\
&=  \min \bigg( \sum_{i=1}^{n_{v_t}} \int_{t=0}^{T} c_1dt + \sum_{i=n_{v_t}+1}^{n} \int_{t=0}^{T} c_0dt \bigg) \nonumber \\
&=  \min (\sum_{i=1}^{n_{v_t}} c_1 T+ \sum_{i=n_{v_t}+1}^{n} c_0 T) \nonumber \\
&=  \min (c_1n_{v_t}T + c_0(n-n_{v_t})T) 
\end{align}
Since by definition $c_0 < c_1$ and $r_s < r_l$, let $ c_0 = r_s$ and $c_1 = r_l$ without loss of generality. Therefore:
\begin{align}
&\phantom{=\ } \min (c_1n_{v_t}T + c_0(n-n_{v_t})T) \nonumber \\
&= \min (r_ln_{v_t} T+ r_s(n-n_{v_t}) T) \nonumber \\
&= \min (\sum_{t=1}^T(r_ln_{v_t} + r_s(n-n_{v_t}))) & \mathllap{[\sum_{i=1}^n x = nx]} \nonumber \\
&\equiv  r_M + \min_{\pi} \mathbb{E} [  \sum_{t=0}^T(r_l n_{v_t} +  r_s (n - n_{v_t})) \mid . ] \nonumber \\
&\equiv Q^{\pi^*}(s_t,a_t) 
\end{align}
Here, the $\equiv$ symbol means that they have a one-to-one mapping or bijection (since $r_M$ is a constant and the reward distribution is fixed). Hence, the optimal cost of navigation within the time horizon $T$ is equivalent to the optimal action-value $Q^{\pi^*}(s_t,a_t)$. Therefore, the optimal policy $\pi^*$ for the given MDP $M$ will generate the optimal paths with minimum costs, satisfying all three constraints of equation (\ref{eq:opt}). $\qed$
\paragraph*{Remark (software versus idealized MDP).}
The proof above analyzes the deterministic MDP in which $\simulate(\cdot)$ follows the nominal dynamics and the agent observes the true state. Our public Gymnasium implementation adds (i) stochastic wind draws and action noise inside the transition, (ii) Gaussian observation noise on the policy inputs, (iii) optional domain randomization of physical parameters between episodes, and (iv) additional dense reward terms that improve empirical learnability. Consequently, finite-sample RL runs should be interpreted as approximating the idealized objective rather than satisfying Theorem~1 symbol-for-symbol. Tightening the theory to the deployed reward and stochastic kernel is left for future work; meanwhile, the idealized MDP clarifies constraint structure and informs reward scaling.

\section{Experiments}
\label{sec:experiments}
\subsection{Robot Dynamics}
Experimental evaluations were conducted for two groups of robots with different dynamics. One group consists of $n$ wheeled robots operating in an environment with obstacles \citep{plaku2021}, and the other group consists of $n$ omnidirectional UAVs operating in a windy environment. The dynamics of the robots are described as follows:
\subsubsection{Wheeled robots with obstacles}  The configuration $(x_i,y_i,\theta_i) \in \config_i$ for robot $i$ has the position $(x_i,y_i)$ and orientation $\theta_i$. The state $(x_i,y_i,\theta_i,\psi_i,v_i) \in \states_i$ includes the steering angle $\psi_i$ and velocity $v_i$. The control $(\alpha, \omega) \in \mathcal{U}_i$ defines the acceleration $\alpha$ and steering rate $\omega$. The motion dynamics are defined as:
\begin{align}
\dot{s} &= f_i(s,u) = (\dot{x}, \dot{y}, \dot{\theta}, \dot{\psi}, \dot{v}), \\
\dot{x} &= v \cos(\theta) \cos(\psi), \quad \dot{y} = v \sin(\theta) \cos(\psi), \\
\dot{\theta} &= v \sin(\psi) / L, \quad  \dot{v} = \alpha, \quad \dot{\psi} = \omega,
\end{align}
where $L$ is the distance from the back to the front wheels. In the software, planar wind drift and actuator noise can be injected so that realized motions deviate from the nominal field shown above.
\subsubsection{UAVs in windy environment} The configuration $(x_i, y_i) \in \config_i$ denotes the center position of robot $i$, and state $(x_i, y_i, v_{x_i}, v_{y_i}) \in \states_i$ denotes the position and velocities along the $x$ and $y$ axis respectively. The control $(\mu_{x_i}, \mu_{y_i}) \in \mathcal{U}_i$ denotes the forces applied along the $x$ and $y$ axis to change its velocity. For the obstacle-free aerial layouts used in training, $\cup_j \config_i \obst_j = \emptyset$ and $\config_i^{\free} = \config_i$. If the wind is blowing at a nominal angle $\beta_a$ and magnitude $v_a$, the \emph{nominal} continuous-time dynamics are
\begin{align}
\dot{s} &= f_i(s,u) = (\dot{x}, \dot{y}, \dot{v_{x}}, \dot{v_{y}}),\\
\dot{v_x} &= \mu_x / M_i , \quad \dot{v_y} = \mu_y / M_i,  \\
\dot{x} &= v_x + v_a \cos(\beta_a), \quad  \dot{y} = v_y + v_a \sin(\beta_a),
\end{align}
where $M_i$ is the mass of robot $i$. At each step the deployed simulator may resample wind magnitude and heading around $(v_a,\beta_a)$, add action noise before integration, and clip speeds for numerical stability.
\subsection{Methodology}
All experiments were performed on a GPU server with  Intel\textsuperscript{\textregistered} Xeon\textsuperscript{\textregistered} Silver 4215 CPU @2.50\,GHz and eight NVIDIA\textsuperscript{\textregistered} A40 GPUs. We crafted \textbf{ten} aerial field configurations and \textbf{eight} wheeled-robot configurations with different field shapes, robot placements, and ROI locations (JSON layouts for UAVs and INI files for wheeled robots in the accompanying repository). We used the Stable-Baselines3 and SB3-Contrib \citep{stable-baselines3} Python libraries for measuring the performance of six deep RL algorithms: Advantage Actor-Critic (A2C) \citep{mnih2016a2c}, Trust Region Policy Optimization (TRPO) \citep{schulman2015trust}, Proximal Policy Optimization (PPO) \citep{schulman2017ppo}, Truncated Mixture of Continuous Distributional Quantile Critics (TQC) \citep{kuznetsov2020tqc}, CrossQ \citep{bhatt2024crossq}, and Augmented Random Search (ARS) \citep{mania2018ars}. We chose both on-policy (A2C, TRPO, PPO, ARS) and off-policy (CrossQ, TQC) RL methods for our experiments. The objective is to compare the learning efficiency of these RL algorithms for different environmental variations in a variety of training conditions. Their performances are measured based on cumulative rewards over each episode.
Unless otherwise noted, each training job uses on the order of $2 \times 10^6$ environment steps per algorithm--layout--seed configuration in the public scripts; shorter budgets may be used during hyperparameter search. Henderson \textit{et al.}\citep{henderson2018deep} emphasized evaluating RL methods on a variety of hyperparameters. Therefore, we performed experiments in two settings: one using the default hyperparameters given in Stable-Baselines3 and SB3-Contrib, and another using hyperparameters selected by the Optuna \citep{optuna_2019} library. Hence, for the two robot groups, we have four primary experiment sets: UAVs with default and tuned hyperparameters, and wheeled robots with default and tuned hyperparameters. We selected both common and model-specific hyperparameters for optimization. For a specific hyperparameter, we used the same range for all algorithms for consistency. The learning rate $\alpha \in [0.0001 - 0.05]$ is used by all algorithms. The generalized advantage estimator $\lambda \in [0.90 - 1.0]$ is used for all algorithms other than ARS and TQC. The value function coefficient \texttt{vf\_coef} $\in [0.2 - 0.7]$, entropy coefficient \texttt{ent\_coef} $\in [0.0 - 0.05]$ and the maximum gradient clipping \texttt{max\_grad\_norm} $\in [0.30 - 0.99]$ are only used in A2C, PPO and CrossQ. The \texttt{buffer\_size} $\in [1000-100000]$ is used on all off-policy algorithms (CrossQ, TQC). We performed 50 different trials in Optuna \citep{optuna_2019} for 0.5 million timesteps. \textbf{Additional switches} in the codebase enable reward ablations, observation ablations, uncertainty ablations, and domain randomization; we summarize the most representative main experiments here and defer exhaustive sweeps to supplementary material.
\subsection{Results and Analysis}
Table~\ref{tab:alg_analysis} shows the results of the four primary experiment sets. The mean and maximum reward for each experiment are highlighted in bold. The range is the absolute difference between the minimum and maximum reward over the timesteps. The average reward threshold for all the experiments for finding the optimal policy is $0.1 \times 10^6$. The threshold is computed as the mean episodic reward averaged over all the environment variations. For UAVs with default hyperparameters, Fig.~\ref{fig:uav_a} shows the performance of each algorithm, where it is evident that CrossQ outperformed all other algorithms by reaching the reward threshold within 500k timesteps. TRPO, PPO, and TQC exhibit similar performances, and ARS appears to be improving gradually. A2C was stuck in a local optimum lower than the others. The maximum reward for TRPO is slightly higher than that of CrossQ, which confirms that CrossQ was also stuck on a local optimum very close to the global optimal reward. For UAVs with random hyperparameters, Fig.~\ref{fig:uav_b} shows the performance of each algorithm, where TRPO and CrossQ achieved comparable maximum rewards. TRPO, PPO, and A2C were robust against changes in hyperparameters, and ARS performed the worst. For wheeled robots with default hyperparameters, Fig.~\ref{fig:wheeled_a} shows the performance of each algorithm. Here, CrossQ was slightly better than ARS, but ARS was more stable and had less variance. A2C had slightly worse performance and was stuck in a local optimum close to the global optimum. For wheeled robots with random hyperparameters, Fig~\ref{fig:wheeled_b} shows that ARS had the best performance, highest stability, and was robust against changes in hyperparameters.
It is well-known that PPO can be strong in highly stochastic benchmarks \citep{schulman2017ppo, henderson2018deep}. Our training runs, however, combine structured dynamics mismatch between platforms with multiplicative wind noise, actuator noise, and (optionally) noisy observations; under this mixture, we still observe competitive TRPO and CrossQ performance in several regimes, suggesting that inductive biases of the algorithms interact nontrivially with the environment's information structure and reward scaling. We caution that episodic returns aggregate dense shaping terms present in the software, so algorithm rankings should be read as empirical sample-complexity comparisons rather than as certificates of optimality for Theorem~1.
\subsection{Simulation}
We utilized the Gymnasium \citep{towers2024gymnasium} Python API with Pygame \citep{pygame} rendering for \emph{both} robot groups during training and evaluation. Additionally, we used CoppeliaSim \citep{coppeliaSim2013} (previously known as V-Rep \citep{vrep2013}) to visualize a 3-D cornfield-style scene in which aerial robots traverse among visually marked target locations; this visualization parallels agricultural patrol routes but should not be confused with the selective \emph{spraying} POMDP studied in our companion paper. Based on the results, the simulations were performed using the trained models for the environment variations 1 and 3 in the scene. Fig.~\ref{fig:sim_both_groups} shows the simulation of environment 1 of both robot groups by running inference on the models trained using the CrossQ \citep{bhatt2024crossq} algorithm. The starting and final positions of the UAVs and target locations (\textcolor{darkgreen}{dark green} markers) in the CoppeliaSim scene are shown in Fig.~\ref{fig:coppelia_set1_i} and  Fig.~\ref{fig:coppelia_set1_f}, respectively, where the \textcolor{red}{red} lines are executed trajectories covering all targets. The starting and final positions of the wheeled robots in Pygame are shown in Fig.~\ref{fig:wheeled_set1_i} and  Fig.~\ref{fig:wheeled_set1_f}, respectively, where the obstacles are shown in \textcolor{gray}{gray} and goal positions are shown in \textcolor{green}{green}.
\subsection{Sensitivity Analysis}
To measure the robustness of the environment design, we ran inference of the UAVs for environment variation 3, where the model was trained using the CrossQ algorithm with no winds ($v_a = 0, \beta_a = 0^{\circ}$), and the inference was run with the wind blowing with variable speeds $v_a$ and angle $\beta_a=30^{\circ}$. Fig. \ref{fig:windy_sim} shows the results where the robots were able to find near-optimal visitation trajectories where $v_a \leq 0.4$. The robustness depends on the maximum and minimum speeds of the UAVs by taking any action. We use $v_{clip} = \pm 5$ units per second for the maximum and minimum speed, which prevents a UAV from gaining unreasonable speed that cannot be visualized in the simulators. Therefore, engineers can use the trained model to generate paths in uncertain and windy environments if $ -\frac{v_{clip}}{10} < v_a < \frac{v_{clip}}{10}$ holds under the tested conditions.
\subsection{Sim-2-Real Analysis}
Although the present study focuses on simulation-based validation, the proposed framework is designed to be extensible to real robotic systems. First, the MDP formulation explicitly incorporates robot dynamics through a generic function $f_i$, enabling direct replacement with system-identified or hardware-calibrated models. Second, environmental disturbances such as wind are modeled within the transition function and varied during inference, serving as a form of structured domain randomization to assess robustness. Additional sensing noise, actuator delays, and model uncertainties can be incorporated directly into the state-transition dynamics without modifying the high-level optimization viewpoint. Furthermore, actuator saturation and finite episode horizons already reflect realistic physical constraints such as battery limits and velocity bounds. Because the environment is formally specified as an MDP, it can also be integrated with reachability analysis, control barrier functions, or temporal-logic-based safety verification prior to deployment. While hardware validation remains future work, these design choices provide a principled pathway for bridging the sim-to-real gap. In software, we additionally provide scripts to quantify performance degradation when injecting observation-level perturbations at deployment time.
\begin{table}[tbp]
\centering
\caption{Experimental results of each algorithm (aggregate statistics; refresh numbers after retraining with ten UAV and eight wheeled layouts).}
\resizebox{\linewidth}{!}{%
\begin{tabular}{
|l|l|
S[table-format=2.4]
S[table-format=2.4]
S[table-format=2.4]
S[table-format=2.4]
S[table-format=2.4]
S[table-format=2.4]|
}
\hline
\multirow{2}{*}{\textbf{Experiment}} & \multirow{2}{*}{\textbf{Param.}} 
& \textbf{A2C} & \textbf{PPO} & \textbf{TRPO} 
& \textbf{ARS} & \textbf{CrossQ} & \textbf{TQC} \\
& & \texttimes ${10^6}$ & \texttimes ${10^6}$ & \texttimes ${10^6}$ & \texttimes ${10^6}$ & \texttimes ${10^6}$ & \texttimes ${10^6}$ \\
\hline
\multirow{4}{1.3cm}{UAV Default}
& Mean  & -0.516027656 & -0.212796252 & -0.095401718 & -0.270643904 & {\bfseries\num{0.142910101}} & -0.197661577 \\
& Max   & -0.100244 & -0.089993898 & {\bfseries\num{0.15937}} & -0.08024 & 0.159349 & 0.132568703 \\
& SD    & 0.243661918 & 0.184584128 & 0.115658208 & 1.483008214 & 0.099024005 & 0.573113877 \\
& Range & 6.609696 & 3.704403102 & 1.67922125 & 29.729134 & 9.961687 & 21.1363947 \\
\Xhline{3\arrayrulewidth}
\multirow{4}{1.3cm}{UAV Random}
& Mean  & -0.243569045 & -0.180676921 & -0.140416732 & -1.534535784 & {\bfseries\num{-0.055438575}} & -0.859217802 \\
& Max   & -0.082030797 & -0.076667 & 0.0067228 & -0.14081 & {\bfseries\num{0.15925}} & 0.0141937 \\
& SD    & 0.09897454 & 0.066854611 & 0.101052405 & 4.378146673 & 0.398611677 & 2.713735699 \\
& Range & 2.356811703 & 0.654988813 & 0.661329175 & 19.761 & 11.761465 & 18.5190737 \\
\Xhline{3\arrayrulewidth}
\multirow{4}{1.3cm}{Wheeled Default}
& Mean  & -0.016639304 & -0.063006647 & -0.04333655 & -0.003461342 & {\bfseries\num{0.084740173}} & -0.067277202 \\
& Max   & 0.019253529 & 0.018938625 & 0.061124766 & 0.019592947 & {\bfseries\num{0.139911266}} & 0.023463559 \\
& SD    & 0.025003901 & 0.036231457 & 0.029397484 & 0.013888309 & 0.083832282 & 0.02349539 \\
& Range & 0.120209717 & 0.122545625 & 0.163848563 & 0.114719588 & 0.240622492 & 0.125124809 \\
\Xhline{3\arrayrulewidth}
\multirow{4}{1.3cm}{Wheeled Random}
& Mean  & -0.095972177 & -0.096296106 & -0.04483161 & {\bfseries\num{-0.00529511}} & -0.083017018 & -0.071808759 \\
& Max   & -0.070700594 & -0.063067281 & 0.016260944 & 0.018952 & {\bfseries\num{0.115718211}} & 0.014393159 \\
& SD    & 0.004539169 & 0.00437403 & 0.027743266 & 0.012253 & 0.01942955 & 0.024753375 \\
& Range & 0.030011867 & 0.038479844 & 0.118246843 & 0.107393719 & 0.216983211 & 0.117600706 \\
\hline
\end{tabular}}
\label{tab:alg_analysis}
\vspace{-10pt}
\end{table}
% \begin{figure}[tbp]
%     \centering
%     \subfloat[UAV Default]{\includegraphics[width=0.499\linewidth]{figures/uav_default_August24.png}\label{fig:uav_a}}
%     \hfil
%     \subfloat[UAV Random]{\includegraphics[width=0.499\linewidth]{figures/uav_random_August24.png}\label{fig:uav_b}}
%     \hfil
%     \subfloat[Wheeled Default]{\includegraphics[width=0.499\linewidth]{figures/wheeled_default_beocat_August24.png}\label{fig:wheeled_a}}
%     \hfil
%     \subfloat[Wheeled Random]{\includegraphics[width=0.499\linewidth]{figures/wheeled_random_beocat_August24.png}\label{fig:wheeled_b}}
%     \caption{Mean episodic reward aggregated across field layouts for the four primary experiment sets: UAVs with default and tuned hyperparameters (a, b), and wheeled robots with default and tuned hyperparameters (c, d).}
%     \label{fig:algorithm_comparison}
%     \vspace{-10pt}
% \end{figure}
% \begin{figure}[tbp]
%     \centering
%     \subfloat[Initial]{\includegraphics[width=0.2\textwidth,height=0.15\textwidth]{figures/env1_coppelia_init.JPG}\label{fig:coppelia_set1_i}}
%     \hfil
%     \subfloat[Final]{\includegraphics[width=0.2\textwidth,height=0.15\textwidth]{figures/env1_coppelia_final.JPG}\label{fig:coppelia_set1_f}} 
%     \hfil
%     \subfloat[Initial]{\includegraphics[width=0.2\textwidth,height=0.15\textwidth]{figures/env1_init.JPG}\label{fig:wheeled_set1_i}}
%     \hfil
%     \subfloat[Final]{\includegraphics[width=0.2\textwidth,height=0.15\textwidth]{figures/env1_final_crossq.JPG}\label{fig:wheeled_set1_f}}    
%     \caption{The initial and final states of the environment 1 for two UAVs in CoppeliaSim (a, b) and for three wheeled robots in Pygame (c, d).}
%     \label{fig:sim_both_groups} 
%     \vspace{-10pt}
% \end{figure}
% \begin{figure*}[tbp]
%     \centering
%     \subfloat[$v_a=0.1$]{\includegraphics[width=0.15\textwidth, height=0.10\textwidth]{figures/env3p_0.1.png}\label{fig:env3p_0.1}} 
%     \hfil
%     \subfloat[$v_a=0.2$]{\includegraphics[width=0.15\textwidth, height=0.10\textwidth]{figures/env3p_0.2.png}\label{fig:env3p_0.2}}
%     \hfil
%     \subfloat[$v_a=0.3$]{\includegraphics[width=0.15\textwidth, height=0.10\textwidth]{figures/env3p_0.3.png}\label{fig:env3p_0.3}}
%     \hfil
%     \subfloat[$v_a=0.4$]{\includegraphics[width=0.15\textwidth, height=0.10\textwidth]{figures/env3p_0.4.png}\label{fig:env3p_0.4}}
%     \hfil
%     \subfloat[$v_a=0.5$]{\includegraphics[width=0.15\textwidth, height=0.10\textwidth]{figures/env3p_0.5.png}\label{fig:env3p_0.5}}
%     \\
%     \subfloat[$v_a=0.1$]{\includegraphics[width=0.15\textwidth, height=0.10\textwidth]{figures/env3s_0.1.jpg}\label{fig:env3s_0.1}} 
%     \hfil
%     \subfloat[$v_a=0.2$]{\includegraphics[width=0.15\textwidth, height=0.10\textwidth]{figures/env3s_0.2.jpg}\label{fig:env3s_0.2}}
%     \hfil
%     \subfloat[$v_a=0.3$]{\includegraphics[width=0.15\textwidth, height=0.10\textwidth]{figures/env3s_0.3.jpg}\label{fig:env3s_0.3}}
%     \hfil
%     \subfloat[$v_a=0.4$]{\includegraphics[width=0.15\textwidth, height=0.10\textwidth]{figures/env3s_0.4.jpg}\label{fig:env3s_0.4}}
%     \hfil
%     \subfloat[$v_a=0.5$]{\includegraphics[width=0.15\textwidth, height=0.10\textwidth]{figures/env3s_0.5.jpg}\label{fig:env3s_0.5}}   
%     \caption{Simulation of environment 3 with winds blowing at an angle $\beta_a=30$° and variable speeds $v_a$ in Pygame (a-e) and in CoppeliaSim (f-j). The trained model was able to find the optimal path for wind speed up to $0.4$ units per second.}
%     \label{fig:windy_sim} 
%     \vspace{-10pt}
% \end{figure*}

\section{Limitations}
\label{sec:limitations}
This work is an improvement of a previous study where the state and action spaces were discrete and finite, and the dynamics were limited to a specific group of robots. The maximum path or episode length (correlated to flight time) for any robot is fixed and enough for every environment variation to visit all regions of interest. But if the episode length (or flight time) is not enough to visit all locations, then we may get a sub-optimal path, for which further experiments need to be performed. All robots in a group are identical and follow the same dynamics. The sample complexity of the state space increases exponentially as the number of robots increases. We tested our environment with up to seven robots. Even though the learned model can reach the reward threshold within a reasonable time, the absolute maximum reward cannot be guaranteed. The hyperparameters were selected based on their performance in environment variation 1 using Optuna, and therefore are not entirely optimal. Collision avoidance is guaranteed for the \emph{idealized} optimal policy of Section~\ref{sec:mdp_theory} through the reward design and formal proof. Learned policies under function approximation, stochastic transitions, and noisy observations may violate constraints at run time; additional safety verification (e.g., model checking via temporal logic) remains important.
Even though the dynamics that we've implemented are relatively simple, it is not a limitation since our approach can be used to implement any complex robot dynamics and variable environments. Despite these limitations, we believe that our work will guide researchers on how to improve the implementation of RL in each development stage, from formally defining the MDP to controlling robots in uncertain and noisy environments.

\section{Conclusion}
\label{sec:conclusion}
We studied learning-based multi-robot navigation in stochastic and uncertain settings while contrasting aerial point-mass dynamics with wheeled bicycle dynamics. An idealized finite-horizon MDP specification and optimization viewpoint clarify constraint structure, while the released Gymnasium software exposes wind, actuator, and sensing noise together with optional domain randomization and ablations for systematic benchmarking. Standard deep RL algorithms train successfully in this setting, and CoppeliaSim visualizations illustrate qualitatively reasonable trajectories under wind perturbations. Future work includes tightening the formal model to match deployed rewards and kernels, hardware validation, and extending the companion spraying formulation to coupled ground--air teams.



%===============================================================================

\clearpage
% The acknowledgments are automatically included only in the final and preprint versions of the paper.
\acknowledgments{If a paper is accepted, the final camera-ready version will (and probably should) include acknowledgments. All acknowledgments go at the end of the paper, including thanks to reviewers who gave useful comments, to colleagues who contributed to the ideas, and to funding agencies and corporate sponsors that provided financial support.}

%===============================================================================

% no \bibliographystyle is required, since the corl style is automatically used.
\bibliography{ref}  % .bib

\end{document}
